%! Equations and Graphs in Both Directions — Unit 2, Lesson 2.3 — Group Activity
\documentclass[10pt]{article}
\usepackage{atda-article}
\usepackage{atda-boxes}
\newcommand{\TallMath}[1]{$\displaystyle #1\rule[-1.4em]{0pt}{3.2em}$}
\setlength{\workrowsep}{6pt}

\begin{document}

\pageheader{Unit 2, Lesson 2.3}{Group Activity --- Both Directions}
% NAMESTRIP: no \namedateperiod here. The student writes their name once, on the
% cover this component is stapled behind. See references/conventions.md.

\vspace{0.15in}

\begin{scenariobox}[The Scenario]{royal}
\small
Half of today's work goes equation $\to$ graph, and half goes graph $\to$ equation. Every time you
write an equation, \textbf{finish with a substitution check}: pick one point you can read off the
graph, put its $x$ into your equation, and confirm you get its $y$. An unchecked equation is a
guess. \textbf{Everyone starts at Tier R.} When your whole group can do Tier R without help, move on.
\end{scenariobox}

\vspace{0.10in}

\boxguard[30]
\begin{tcolorbox}[colback=white, colframe=black!40, title=\textbf{Tier R --- Remediate},
    fonttitle=\bfseries, arc=2mm, left=3mm, right=3mm, top=2mm, bottom=2mm, breakable]
\small
\textbf{Equation $\to$ graph.} Build the graph of $y=(x+1)^2-4$ from the parent.

\begin{enumerate}[leftmargin=1.6em, itemsep=4pt, topsep=2pt]
  \item Name the parent, then give the anchor point of the image. Say which slide came from the
        number inside the parentheses and which came from the number outside.

\writelines{2}

  \item Complete the table, then plot your five points on the grid and join them with a smooth
        curve. The parent is already drawn.
\smallskip

\renewcommand{\arraystretch}{1.3}
\begin{tabular}{@{}|c|c|c|c|c|c|@{}}
\hline
$x$ & $-3$ & $-2$ & $-1$ & $0$ & $1$ \\ \hline
$y=(x+1)^2-4$ & \blank{1.0cm} & \blank{1.0cm} & \blank{1.0cm} & \blank{1.0cm} & \blank{1.0cm} \\
\hline
\end{tabular}
\smallskip

\begin{center}
\begin{tikzpicture}
\begin{axis}[
    width=9.4cm, height=5.4cm,
    xlabel={$x$}, ylabel={$y$},
    xmin=-4.4, xmax=3.4, ymin=-5.4, ymax=5.4,
    xtick={-4,-3,-2,-1,0,1,2,3}, ytick={-5,-4,-3,-2,-1,0,1,2,3,4,5},
    grid=both, grid style={linegray!45}, axis lines=left,
    tick label style={font=\tiny}, label style={font=\scriptsize},
    legend entries={{$y=x^2$ (the parent)}},
    legend style={font=\scriptsize, draw=none, fill=none,
      at={(0.5,1.02)}, anchor=south, legend columns=2}]
  \addplot[royal, very thick, domain=-2.2:2.2, samples=80]{x^2};
\end{axis}
\end{tikzpicture}
\end{center}

  \item \textbf{Graph $\to$ equation, the easy direction.} A parabola has exactly the same shape as
        $y=x^2$, and its lowest point is at $(4,0)$. Write its equation, then check it by
        substituting $x=6$.

\writelines{2}

  \item Which of the domain and the range did the $-4$ in $y=(x+1)^2-4$ move? Give the range of the
        image in interval notation.

\writelines{1}
\end{enumerate}
\end{tcolorbox}

\vspace{0.10in}

\boxguard[30]
\begin{tcolorbox}[colback=white, colframe=black!40,
    title=\textbf{Tier A --- Approaching Proficiency},
    fonttitle=\bfseries, arc=2mm, left=3mm, right=3mm, top=2mm, bottom=2mm, breakable]
\small
\textbf{Graph $\to$ equation.} Now the pictures are handed to you and the shapes are not all the
parent's.

\begin{center}
\begin{tabular}{@{}c@{\hspace{0.6cm}}c@{}}
\begin{tikzpicture}
\begin{axis}[
    width=6.6cm, height=4.4cm, title={\scriptsize (a) a parabola},
    xmin=0, xmax=6.4, ymin=-4.4, ymax=11.4,
    xtick={0,1,2,3,4,5,6}, ytick={-3,0,3,6,9}, minor y tick num=2,
    grid=both, grid style={linegray!45}, minor grid style={linegray!30},
    axis lines=left,
    tick label style={font=\tiny}]
  \addplot[royal, very thick, domain=1:5, samples=90]{3*(x-3)^2-2};
\end{axis}
\end{tikzpicture}
&
\begin{tikzpicture}
\begin{axis}[
    width=6.6cm, height=4.4cm, title={\scriptsize (b) an exponential},
    xmin=-3.4, xmax=2.4, ymin=0, ymax=7.4,
    xtick={-3,-2,-1,0,1,2}, ytick={0,2,4,6}, minor y tick num=1,
    grid=both, grid style={linegray!45}, minor grid style={linegray!30},
    axis lines=left,
    tick label style={font=\tiny}]
  \addplot[linegray, thick, dashed, domain=-3.3:2.3, samples=2]{2};
  \addplot[royal, very thick, domain=-3.3:2.2, samples=90]{2^x+2};
\end{axis}
\end{tikzpicture}
\end{tabular}
\end{center}

\begin{enumerate}[leftmargin=1.6em, itemsep=4pt, topsep=2pt]
  \item \textbf{Graph (a).} Read the vertex, then use the point $(4,1)$ to find $a$. Write the
        equation in transformation form.
\begin{work}
  1 &= a(4-3)^2-2 \\
  1 &= a-2 \\
  3 &= a
\end{work}

  \item Check your equation from item 1 at $x=5$, and say what the graph must show there if you are
        right.

\writelines{2}

  \item \textbf{Graph (b).} The dashed line is the flat line this curve crawls along. Give its
        equation, give the point the curve passes through on the $y$-axis, and write the function.
        Then state the range in interval notation --- bracket or parenthesis, and why.

\writelines{3}

  \item \textbf{In context.} A gym charges a joining fee plus a flat monthly rate. The graph of the
        total cost $C(x)$ after $x$ months is a straight line through $(0,40)$ and $(4,100)$.
        \begin{enumerate}[label=(\alph*), leftmargin=1.6em, itemsep=3pt, topsep=3pt]
          \item Find the monthly rate from the two points, and write $C(x)$.

\writelines{1}

          \item Use your equation to find the total cost after $7$ months.
\begin{work}
  C(7) &= 15(7)+40 \\
       &= 105+40 \\
       &= 145
\end{work}

          \item The notes said that for a \emph{line} you report two points rather than naming a
                transformation. Give the reason, in your own words.

\writelines{2}
        \end{enumerate}
\end{enumerate}
\end{tcolorbox}

\vspace{0.10in}

\boxguard[30]
\begin{tcolorbox}[colback=white, colframe=black!40, title=\textbf{Tier E --- Extension},
    fonttitle=\bfseries, arc=2mm, left=3mm, right=3mm, top=2mm, bottom=2mm, breakable]
\small

\begin{enumerate}[leftmargin=1.6em, itemsep=4pt, topsep=2pt]
  \item \textbf{Find the typo.} Each student below typed something into a graphing calculator and
        got a graph that does not match what they were asked for. Name the error and give the
        keystrokes that fix it.
        \begin{enumerate}[label=(\alph*), leftmargin=1.6em, itemsep=3pt, topsep=3pt]
          \item Asked for: the exponential parent moved \emph{right} $5$. Typed:
                \texttt{Y1=2\^{}X-5}.

\writelines{2}

          \item Asked for: the quadratic parent \emph{reflected over the $x$-axis}. Typed:
                \texttt{Y1=(-X)\^{}2}. The screen showed a parabola opening \emph{upward}.

\writelines{2}
        \end{enumerate}

  \item \textbf{The window.} A classmate typed $y=(x-15)^2-40$ correctly and pressed
        \texttt{ZOOM}\,\texttt{6}, which sets the standard window $-10 \le x \le 10$,
        $-10 \le y \le 10$. The screen came up empty.
        \begin{enumerate}[label=(\alph*), leftmargin=1.6em, itemsep=3pt, topsep=3pt]
          \item Explain why the screen is empty, using the anchor point.

\writelines{1}

          \item Give four window values that would show the vertex and both sides of the curve, and
                say how you chose them \emph{from the equation}, without graphing anything.

\writelines{2}
        \end{enumerate}

  \item \textbf{Justify.} A classmate says: ``My calculator screen looks exactly like the printed
        graph, so my equation must be right.'' Explain why that is not proof --- use $y=(x-3)^2$ and
        $y=(x-3)^2+0.5$ as your example --- and then say what \emph{would} be proof.

\writelines{3}
\end{enumerate}
\end{tcolorbox}

\end{document}
